\chapter{Capas de cizalla}
\label{sec:shearing-layers}

\section{La metáfora}

Stewart Brand, en \emph{How Buildings Learn}, observa que un edificio
no cambia a la misma velocidad en todas sus partes: los cimientos
duran décadas, la estructura medio siglo, la piel del edificio
treinta años, los servicios (plomería, electricidad) una década,
la distribución interior unos cinco años, y las cosas que ponemos
encima de las mesas cambian a diario. Cada capa cede (\emph{shears})
a distinta velocidad; si se las acopla rígidamente, el ritmo rápido
arrastra al lento y se termina rompiendo.

GraphHunter aplica la misma idea al software. Hay piezas que
cambian poco (algoritmos de grafo, intersección AVX2, operaciones
FFI) y piezas que cambian seguido (un nuevo parser de log, una
nueva tool MCP, una pantalla nueva del Tauri). Querer aplicar
exactamente los mismos estándares de rigor a ambas mata la
velocidad; querer aplicar los mismos estándares de flexibilidad
mata la corrección.

\section{Tres capas}

El repositorio se organiza en tres capas:

\begin{description}
    \item[Núcleo (\filepath{core/})] — Algoritmos y tipos
    fundamentales. Contratos matemáticos explícitos, property-based
    testing, benchmarks. Abstracciones sólo cuando hay variación
    real. Cambia lento; cada cambio pesa.
    \item[Plataforma (\filepath{platform/})] — Superficies de
    extensión: API, DSL, parsers, fuentes de log, canonicalización,
    MCP. SOLID clásico aplica aquí: OCP e ISP, traits como contratos
    públicos. Cambia a ritmo medio; lo que se agrega debe mantener
    compatibilidad.
    \item[Aplicaciones (\filepath{apps/}, \filepath{platform/mcp/})]
    — Piel visible: Tauri (GUI), CLI, MCP server TypeScript. YAGNI
    agresivo, código pragmático y directo. Cambia rápido; lo
    importante es que el cambio sea obvio.
\end{description}

La regla de acoplamiento entre capas es \textbf{unidireccional hacia
abajo}: \filepath{apps/} puede depender de \filepath{platform/}, y
\filepath{platform/} puede depender de \filepath{core/}. Lo inverso
no se permite. Esto garantiza que un cambio en una pantalla del
Tauri nunca force un cambio en el algoritmo de matching, y que un
nuevo parser no pueda corromper el núcleo algorítmico.

\section{Criterios de clasificación}

La decisión de qué va en qué capa no es una votación. Se aplica,
en orden, este checklist:

\begin{enumerate}
    \item ¿El módulo implementa un \textbf{algoritmo con contrato
    matemático} (DFS, intersección, scoring)? Núcleo.
    \item ¿El módulo define un \textbf{trait de extensión} que
    terceros (o futuros módulos) van a implementar (\code{LogParser},
    \code{DslExtension}, \code{Tool})? Plataforma.
    \item ¿El módulo es una \textbf{superficie pública versionable}
    (endpoints HTTP, tools MCP, DTOs, C ABI, DSL)? Plataforma.
    \item ¿El módulo es \textbf{pass-through} o binding 1:1 a una
    pieza ya existente (Tauri command que reenvía a la API, botón
    que llama a un endpoint)? Aplicación.
    \item ¿El módulo es \textbf{setup} de proceso (main, argparse,
    router de HTTP, IPC bridge)? Aplicación.
\end{enumerate}

Cuando un módulo cumple varios criterios, el que aparezca primero
gana (algoritmo $>$ trait $>$ superficie $>$ pass-through $>$ setup).

\section{Diferencia con SOLID uniforme}

Un enfoque SOLID uniforme tratar\'ia a todas las piezas del código
con la misma vara: interfaces para todo, inversión de dependencias
para todo, apertura/cerrada para todo. El costo es real: traits con
una sola implementación que sólo existen para hacer ``tests
unitarios'' con mocks, arquitecturas hexagonales con tres niveles
de indirección entre un HTTP handler y una query SQL, y una
cognitive load desproporcionada para cambios triviales.

Shearing layers no niega SOLID; lo \emph{localiza}. SOLID reina en
\filepath{platform/}, donde la extensibilidad es una feature de
primera clase. En \filepath{core/} se prefiere tipos fuertes,
invariantes codificadas y property tests; las abstracciones nacen
sólo cuando hay una segunda implementación real. En \filepath{apps/}
se prefiere código imperativo plano, legible sin
saltos de archivo.

\section{Cómo se refleja en el árbol}

\begin{lstlisting}[language=, basicstyle=\ttfamily\small]
GraphHunter/
|-- core/
|   |-- graph-engine/       # algoritmos de grafo, DFS, anomaly
|   |-- matcher-ffi/        # fachada segura hacia libgraphmatch
|   |-- gnn/                # scorer ONNX (v1 estable, v2 gate)
|   `-- gnn/src/v2/         # experimental, feature-gated
|-- platform/
|   |-- api/                # GraphHunterApi (facade)
|   |-- canonical/          # OCSF v1.4 projection
|   |-- constrained-decode/ # guardrail LLM
|   |-- dsl/                # hypothesis DSL
|   |-- local-llm/          # trait LocalLlm + backends
|   |-- parsers/            # trait LogParser + 7 impl
|   |-- siem/               # Sentinel/Elastic streaming
|   |-- sources/            # trait LogSource
|   |-- vrl/                # FieldConfig -> VRL compiler
|   `-- mcp/                # TypeScript MCP server (54 tools)
|-- apps/
|   |-- tauri/              # GUI Tauri + React
|   `-- cli/                # CLI (usa platform/api in-proc)
|-- libgraphmatch/          # C++20 AVX2, no Rust
`-- legacy/                 # gateway Go, archivado
\end{lstlisting}

\section{Lo que ganamos}

Cuatro propiedades se siguen directamente de esta estructura:

\begin{itemize}
    \item \textbf{Cambios aislados}: reemplazar un parser no afecta
    al núcleo; cambiar el scoring no afecta al DSL; reorganizar la
    GUI no afecta a nada del backend.
    \item \textbf{Revisión proporcional}: un PR que toca sólo
    \filepath{apps/} pasa con revisión ligera; uno que toca
    \filepath{core/} exige benchmarks y pruebas de propiedad; uno que
    toca \filepath{platform/} exige compatibilidad de API.
    \item \textbf{Tests apropiados}: property-based y equivalencia
    en el núcleo, contractuales en la plataforma, smoke en
    aplicaciones. Sin uniformar a lo más costoso.
    \item \textbf{Onboarding escalonado}: un contribuidor nuevo
    empieza tocando \filepath{apps/}, gana contexto, después migra a
    \filepath{platform/}, y sólo entonces toca \filepath{core/}.
\end{itemize}

El resto del libro asume este mapa. El capítulo siguiente muestra
el grafo de dependencias que lo hace verificable.
